\documentclass[11pt]{article}

\usepackage[margin=1in]{geometry}
\usepackage{amsmath,amssymb,amsthm}
\usepackage{mathtools}
\usepackage{booktabs}
\usepackage{enumitem}
\usepackage[colorlinks=true,linkcolor=blue,citecolor=blue]{hyperref}

\newtheorem{theorem}{Theorem}
\newtheorem{lemma}{Lemma}
\newtheorem{proposition}{Proposition}
\newtheorem{corollary}{Corollary}
\theoremstyle{definition}
\newtheorem{remark}{Remark}

\DeclareMathOperator*{\argmin}{arg\,min}
\newcommand{\R}{\mathbb{R}}
\newcommand{\Ln}{L_n}
\newcommand{\thetast}{\theta^{*}}
\newcommand{\Mbar}{\overline{M}}
\newcommand{\Mperp}{M^{\perp}}
\newcommand{\Mbarperp}{\overline{M}^{\perp}}
\newcommand{\Phid}{\Phi^{*}}
\newcommand{\PsM}{\Psi(\Mbar)}
\newcommand{\norm}[1]{\lVert #1 \rVert}
\newcommand{\ip}[2]{\langle #1,\, #2 \rangle}
\newcommand{\Fcal}{\mathcal{F}}
\newcommand{\Ccal}{\mathcal{C}}

\title{\bf On the Estimation--Error Constants in Wainwright,\\
\emph{High-Dimensional Statistics}, Chapter 9:\\[2pt]
\large A Lean 4 Formalization Report}
\author{StatLean formalization notes (Batch 6)}
\date{June 2026}

\begin{document}
\maketitle

\begin{abstract}
While formalizing Chapter~9 of Wainwright's \emph{High-Dimensional Statistics} (regularized
$M$-estimators, decomposable regularizers, restricted strong convexity) in Lean~4, we found that
the leading \emph{constants} of the master estimation-error bound (Theorem~9.19 / Corollary~9.20)
and of its generalized-linear-model specialization (Corollaries~9.26--9.27) are not recoverable
verbatim by a fully rigorous, finite-sample proof. The book's squared-error constant $9$ becomes a
provable $144$; the $\ell_1$ constant becomes $48$; the restricted-strong-convexity slack tightens
from $\tau^2\Psi^2\le\kappa/64$ to $\tau^2\Psi^2\le\kappa/128$; and Corollary~9.27 must be stated with
the \emph{same} regularization strength $\lambda=4BC\{\sqrt{\log d/n}+\delta\}$ as Corollary~9.26,
not the smaller $2BC\{\cdots\}$. We pinpoint the mechanism --- the restricted-strong-convexity
\emph{defect} term is quadratic in the error on the constraint cone and therefore competes directly
with the curvature term, rather than being a negligible lower-order perturbation --- exhibit the
explicit slack budget that yields the provable constants, and verify that the statistical
\emph{rate} $O\!\big(s\lambda^2/\kappa^2\big)$ is unchanged. All deviations are recorded in the
corresponding Lean docstrings, per the project convention that textbook constants are stated only up
to small factors.
\end{abstract}

\section{Context and summary of findings}

The Lean development \texttt{StatLean/HighDimensionalStatistics/MEstimator/} formalizes the general
$M$-estimation theory of Wainwright, Chapter~9. Every theorem is machine-checked and
\texttt{\#print axioms}--clean (no \texttt{sorryAx}). A formal proof, unlike a textbook proof, cannot
absorb a quantity into ``$o(1)$'' or carry an implicit ``for $n$ large enough'': every constant must be
\emph{exhibited} and every inequality discharged at finite $n$. This discipline surfaced a family of
constant discrepancies, summarized in Table~\ref{tab:summary}.

\begin{table}[h]
\centering
\renewcommand{\arraystretch}{1.25}
\begin{tabular}{@{}lll@{}}
\toprule
Quantity & Wainwright (book) & Provable (Lean) \\
\midrule
Thm 9.19 / Cor 9.20 \quad $\norm{\widehat\theta-\thetast}^2 \le (\,\cdot\,)\,(\lambda/\kappa)^2\PsM^2$
 & $9$ & $144$ \\
Cor 9.20 / 9.26 \quad $\Phi(\widehat\theta-\thetast) \le (\,\cdot\,)\,(\lambda/\kappa)\PsM^2$
 & $\approx 6$ & $48$ \\
RSC slack (Thm 9.19) & $\tau_n^2\PsM^2 \le \kappa/64$ (effective) & $\tau_n^2\PsM^2 \le \kappa/128$ \\
Curvature slack (Thm 9.24) & $\tau_n\PsM^2 < \kappa/32$ & $\kappa/32$ \;(matches) \\
$\ell_\infty$ bound (Thm 9.24) & $\Phid(\widehat\theta-\thetast)\le 3\lambda/\kappa$ & $3\lambda/\kappa$ \;(matches) \\
$\lambda$ (Cor 9.26) & $4BC\{\sqrt{\log d/n}+\delta\}$ & matches \\
$\lambda$ (Cor 9.27) & $2BC\{\sqrt{\log d/n}+\delta\}$ & \textbf{must be} $4BC\{\cdots\}$ \\
Good-event probability & $1-2e^{-2n\delta^2}$ & matches \\
\bottomrule
\end{tabular}
\caption{Book vs.\ provable constants. The statistical \emph{rate} is identical throughout; only the
leading constants and the (quantitative) slack thresholds change.}
\label{tab:summary}
\end{table}

The bulk of this note explains the first row --- the $9\to144$ inflation of the squared-error
constant --- which is the deepest of the discrepancies and the source of the others. Sections~\ref{sec:ell1}
and~\ref{sec:linf} treat the $\ell_1$ constant and the Corollary~9.27 regularization-strength issue.

\section{Setup and the master inequality}\label{sec:setup}

Let $E$ be a finite-dimensional real inner-product space, $\Ln:E\to\R$ a convex, differentiable
empirical cost, and $\Phi$ a regularizer (a seminorm) with dual $\Phid$. The estimator is
\[
  \widehat\theta \in \argmin_{\theta\in E}\;\big\{\Ln(\theta) + \lambda\,\Phi(\theta)\big\},
  \qquad \widehat\Delta := \widehat\theta-\thetast .
\]
Following Wainwright we work with a subspace pair $M\subseteq\Mbar$ (the model and its companion),
orthogonal projections $\Delta_{\Mbar},\Delta_{\Mbarperp}$ with respect to $\ip{\cdot}{\cdot}$, and:

\begin{itemize}[leftmargin=1.4em]
\item \textbf{Restricted strong convexity (Def.\ 9.15).} For $\norm{\Delta}\le R$,
\begin{equation}\label{eq:rsc}
  \underbrace{\Ln(\thetast+\Delta)-\Ln(\thetast)-\ip{\nabla\Ln(\thetast)}{\Delta}}_{=:~\mathcal{E}_n(\Delta)}
  \;\ge\; \tfrac{\kappa}{2}\norm{\Delta}^2 \;-\; \tau_n^2\,\Phi^2(\Delta).
\end{equation}
The term $\tau_n^2\Phi^2(\Delta)$ is the \emph{tolerance} or \emph{defect}: it is what prevents
$\Ln$ from being globally strongly convex in high dimensions.
\item \textbf{Good event (eq.\ 9.46).} $\mathfrak{G}(\lambda)=\{\Phid(\nabla\Ln(\thetast))\le\lambda/2\}$.
\item \textbf{Decomposability (Def.\ 9.9)} $\Rightarrow$ \textbf{cone membership (Prop.\ 9.13):} on
$\mathfrak{G}(\lambda)$,
\begin{equation}\label{eq:cone}
  \Phi(\widehat\Delta_{\Mbarperp}) \;\le\; 3\,\Phi(\widehat\Delta_{\Mbar}) + 4\,\Phi(\thetast_{\Mperp}).
\end{equation}
\item \textbf{Subspace Lipschitz constant (Def.\ 9.18).} $\Phi(u_{\Mbar})\le\PsM\,\norm{u}$.
\end{itemize}

Define the optimization-theoretic function
\begin{equation}\label{eq:Fcal}
  \Fcal(\Delta) := \Ln(\thetast+\Delta)-\Ln(\thetast) + \lambda\big[\Phi(\thetast+\Delta)-\Phi(\thetast)\big].
\end{equation}
Optimality of $\widehat\theta$ gives $\Fcal(\widehat\Delta)\le0$.

\begin{lemma}[Sphere positivity, Lemma 9.21]\label{lem:sphere}
$\Fcal$ is convex with $\Fcal(0)=0$. If $\Fcal(\Delta)>0$ for every $\Delta$ with $\norm{\Delta}=r$,
then $\norm{\widehat\Delta}< r$.
\end{lemma}
\begin{proof}[Idea]
If $\norm{\widehat\Delta}\ge r$, the point $\widetilde\Delta=(r/\norm{\widehat\Delta})\,\widehat\Delta$
lies on the sphere, and convexity with $\Fcal(0)=0$ gives
$\Fcal(\widetilde\Delta)\le (r/\norm{\widehat\Delta})\,\Fcal(\widehat\Delta)\le0$, contradicting
positivity on the sphere.
\end{proof}

Thus the whole game is to find the smallest radius $r=\varepsilon_n$ above which $\Fcal>0$ on the
sphere; then $\norm{\widehat\Delta}^2\le\varepsilon_n^2$.

\section{The lower bound on $\Fcal$}\label{sec:lower}

We specialize to the case $\thetast\in M$ (Corollary~9.20, and the GLM Corollary~9.26), where
$\Phi(\thetast_{\Mperp})=0$, so the cone~\eqref{eq:cone} becomes
$\Phi(\Delta_{\Mbarperp})\le 3\Phi(\Delta_{\Mbar})$ and hence
\begin{equation}\label{eq:Phicone}
  \Phi(\Delta) \le \Phi(\Delta_{\Mbar})+\Phi(\Delta_{\Mbarperp}) \le 4\,\Phi(\Delta_{\Mbar}) \le 4\,\PsM\,\norm{\Delta}.
\end{equation}
Combining~\eqref{eq:rsc}, the good event $\ip{\nabla\Ln(\thetast)}{\Delta}\ge-\tfrac{\lambda}{2}\Phi(\Delta)$,
the decomposability lower bound
$\Phi(\thetast+\Delta)-\Phi(\thetast)\ge\Phi(\Delta_{\Mbarperp})-\Phi(\Delta_{\Mbar})$, and
$\Phi(\Delta)\le\Phi(\Delta_{\Mbar})+\Phi(\Delta_{\Mbarperp})$, one obtains after cancellation
\begin{equation}\label{eq:star}
  \Fcal(\Delta) \;\ge\; \tfrac{\kappa}{2}\norm{\Delta}^2 \;-\; \tau_n^2\,\Phi^2(\Delta)
                 \;-\; \tfrac{3\lambda}{2}\,\Phi(\Delta_{\Mbar})
                 \;+\; \tfrac{\lambda}{2}\,\Phi(\Delta_{\Mbarperp}).
\end{equation}
Dropping the nonnegative last term and using $\Phi(\Delta_{\Mbar})\le\PsM\norm{\Delta}$,
\begin{equation}\label{eq:starstar}
  \boxed{\;\Fcal(\Delta) \;\ge\; \tfrac{\kappa}{2}\norm{\Delta}^2 \;-\; \tau_n^2\,\Phi^2(\Delta)
                 \;-\; \tfrac{3\lambda}{2}\,\PsM\,\norm{\Delta}.\;}
\end{equation}
Inequality~\eqref{eq:starstar} is the crux: \emph{the constant is decided entirely by how the middle
defect term $\tau_n^2\Phi^2(\Delta)$ is handled.}

\section{Why the book's constant $9$ is not formally recoverable}\label{sec:why}

\paragraph{The book's reading.} Suppose one treats the defect $\tau_n^2\Phi^2(\Delta)$ as
\emph{negligible} relative to the curvature term --- i.e.\ one drops it, or argues that under
``$\tau_n^2\PsM^2/\kappa\to0$'' it is asymptotically lower order. Then \eqref{eq:starstar} reads
\[
  \Fcal(\Delta) \;\gtrsim\; \tfrac{\kappa}{2}\norm{\Delta}^2 - \tfrac{3\lambda}{2}\PsM\norm{\Delta}
  \;=\; \norm{\Delta}\Big(\tfrac{\kappa}{2}\norm{\Delta} - \tfrac{3\lambda}{2}\PsM\Big),
\]
which is positive once $\norm{\Delta} > 3\lambda\PsM/\kappa$. By Lemma~\ref{lem:sphere},
$\norm{\widehat\Delta}^2 \le 9\,(\lambda/\kappa)^2\PsM^2$. This is exactly Wainwright's
eq.~(9.49b): the constant $9$ is the square of the curvature-to-slope ratio
$\tfrac{3\lambda}{2}\PsM\big/\tfrac{\kappa}{2}=3\lambda\PsM/\kappa$ with the \emph{full} coefficient
$\kappa/2$ surviving in front of $\norm{\Delta}^2$.

\paragraph{Why this is not a finite-sample proof.} The defect term is \emph{not} lower order at the
critical radius. On the cone, \eqref{eq:Phicone} gives $\Phi(\Delta)\le 4\PsM\norm{\Delta}$, so
\begin{equation}\label{eq:defectquad}
  \tau_n^2\,\Phi^2(\Delta) \;\le\; 16\,\tau_n^2\,\PsM^2\,\norm{\Delta}^2 ,
\end{equation}
which is \emph{quadratic} in $\norm{\Delta}$ --- the same order as the curvature term it is being
subtracted from. There is no radius at which it is negligible relative to $\tfrac{\kappa}{2}\norm{\Delta}^2$;
their ratio is the constant $32\tau_n^2\PsM^2/\kappa$, independent of $\norm{\Delta}$. Consequently the
defect \emph{permanently degrades the effective curvature coefficient}:
\begin{equation}\label{eq:ceff}
  \Fcal(\Delta) \;\ge\; \Big(\tfrac{\kappa}{2} - c_{\mathrm{def}}\,\tau_n^2\PsM^2\Big)\norm{\Delta}^2
                 - \tfrac{3\lambda}{2}\PsM\norm{\Delta},
\end{equation}
and the leading constant is governed by $c := \tfrac{\kappa}{2}-c_{\mathrm{def}}\tau_n^2\PsM^2$, which is
\emph{strictly less than} $\kappa/2$ for any $\tau_n>0$. The book's $9$ corresponds to the unattainable
idealization $c=\kappa/2$.

In the formalized \emph{general} theorem ($M\subseteq\Mbar$, $\thetast\notin M$ allowed) the
approximation cross-term $4\Phi(\thetast_{\Mperp})$ in the cone~\eqref{eq:cone} forces the splitting
$(a+b)^2\le 2a^2+2b^2$ when bounding $\Phi^2(\Delta)$, doubling the defect coefficient to
$c_{\mathrm{def}}=32$. Corollary~9.20 is then obtained by \emph{specializing} this single general bound
to $\thetast\in M$, so it \emph{inherits} the general constant rather than re-deriving a tighter one.
This is the precise Lean statement (file \texttt{Bound.lean}):
\[
  c \;=\; \tfrac{\kappa}{2} - 32\,\tau_n^2\,\PsM^2 .
\]

\section{The explicit slack budget and the provable constant $144$}\label{sec:fix}

To turn~\eqref{eq:ceff} into a finite-sample theorem we must commit to a \emph{quantitative} slack and
then \emph{budget} the surviving curvature across the competing terms. The choice made (and verified) is
\begin{equation}\label{eq:slack}
  \tau_n^2\,\PsM^2 \;\le\; \frac{\kappa}{128}
  \quad\Longrightarrow\quad
  c \;=\; \tfrac{\kappa}{2}-32\,\tau_n^2\PsM^2 \;\ge\; \tfrac{\kappa}{2}-32\cdot\tfrac{\kappa}{128}
        \;=\; \tfrac{\kappa}{4}.
\end{equation}
(The book's implicit threshold, read off the same computation with $c_{\mathrm{def}}=16$ and demanding
only $c\ge\kappa/4$, is $\tau_n^2\PsM^2\le\kappa/64$; the general-case factor of two tightens it to
$\kappa/128$.) The surviving budget $\kappa/4$ must simultaneously dominate two competing quantities ---
the linear slope term and the approximation-error term of $\varepsilon_n^2$ --- and reserve a margin.
The Lean proof (\texttt{fcal\_pos\_on\_sphere}) makes the split explicit:
\begin{equation}\label{eq:split}
  \frac{\kappa}{4}\;=\;\underbrace{\frac{\kappa}{8}}_{\text{vs.\ }\frac{3\lambda}{2}\PsM\norm{\Delta}}
                   \;+\;\underbrace{\frac{\kappa}{16}}_{\text{approx.\ error }2\lambda\Phi_0+32\tau_n^2\Phi_0^2}
                   \;+\;\underbrace{\frac{\kappa}{16}}_{\text{margin}},
  \qquad \Phi_0:=\Phi(\thetast_{\Mbarperp}).
\end{equation}
The decisive sub-budget is the first: requiring $\tfrac{3\lambda}{2}\PsM\norm{\Delta}\le\tfrac{\kappa}{8}\norm{\Delta}^2$
forces
\begin{equation}\label{eq:radius}
  \norm{\Delta} \;\ge\; \frac{\tfrac{3\lambda}{2}\PsM}{\kappa/8} \;=\; \frac{12\,\lambda\,\PsM}{\kappa},
  \qquad\text{hence}\qquad
  \norm{\widehat\Delta}^2 \;\le\; 144\,\frac{\lambda^2}{\kappa^2}\,\PsM^2 .
\end{equation}
Thus the provable squared-error constant is $144=12^2$, where $12 = \tfrac{3}{2}\big/\tfrac{1}{8}$
replaces the book's $3 = \tfrac{3}{2}\big/\tfrac{1}{2}$: the effective curvature coefficient has fallen
from $\kappa/2$ to $\kappa/8$, a factor of $4$ on the radius and $16$ on the squared error,
$9\cdot16=144$. The full $\varepsilon_n^2$ (Lean \texttt{epsilonSq}) is correspondingly
\[
  \varepsilon_n^2 \;=\; 144\,\frac{\lambda^2}{\kappa^2}\,\PsM^2
     \;+\; \frac{32}{\kappa}\Big(\lambda\,\Phi_0 + 16\,\tau_n^2\,\Phi_0^2\Big),
\]
to be compared with Wainwright's eq.~(9.47), $9(\lambda/\kappa)^2\PsM^2+\tfrac{8}{\kappa}\{\lambda\Phi_0+16\tau_n^2\Phi_0^2\}$.

\section{The $\ell_1$ / regularizer constant ($6\to48$)}\label{sec:ell1}

The regularizer-error bound (eq.\ 9.48a) follows from the $\ell_2$ bound and the cone. With
$\thetast\in M$, \eqref{eq:Phicone} and~\eqref{eq:radius} give
\[
  \Phi(\widehat\Delta) \;\le\; 4\,\PsM\,\norm{\widehat\Delta}
     \;\le\; 4\,\PsM\cdot\frac{12\,\lambda\,\PsM}{\kappa}
     \;=\; 48\,\frac{\lambda}{\kappa}\,\PsM^2 .
\]
For the GLM Lasso, $\PsM^2=\Psi(M(S))^2=s$ (Lean \texttt{subspaceLip\_l1\_suppSubmodule}), giving
$\norm{\widehat\Delta}_1\le 48\,s\lambda/\kappa$. The book's constant is smaller (of order $6$); the
inflation factor is exactly the radius factor $4$ carried through the relation $\Phi(\widehat\Delta)\le
4\PsM\norm{\widehat\Delta}$, i.e.\ $48 = 4\times 12$, mirroring $144=12^2$.

\section{Corollary 9.27: the regularization-strength inconsistency}\label{sec:linf}

The $\ell_\infty$ rate (Corollary~9.27) is sometimes quoted with the smaller tuning
$\lambda = 2BC\{\sqrt{\log d/n}+\delta\}$, the constant that governs the \emph{deterministic} dual-norm
bound (Theorem~9.24, $\Phid(\widehat\theta-\thetast)\le 3\lambda/\kappa$). This is inconsistent with the
\emph{probabilistic} content of the corollary. The high-probability guarantee rests on the good event
$\mathfrak{G}(\lambda)=\{\norm{\nabla\Ln(\thetast)}_\infty\le\lambda/2\}$, and the union-bound estimate
(Lean \texttt{good\_event\_highProb}) shows
\[
  \Pr\!\big(\mathfrak{G}(\lambda)\big) \ge 1-2e^{-2n\delta^2}
  \qquad\text{requires}\qquad
  \lambda \;\ge\; 4BC\Big\{\sqrt{\tfrac{\log d}{n}}+\delta\Big\}.
\]
Indeed, choosing $t=\lambda/2$ in the coordinatewise sub-Gaussian tail
$\Pr(\norm{\nabla\Ln(\thetast)}_\infty>t)\le 2d\,e^{-t^2/(2B^2C^2/n)}$ and demanding the right-hand side
collapse to $2e^{-2n\delta^2}$ uses $\big(\tfrac{\lambda}{2}\big)^2\ge 4B^2C^2(\log d/n+\delta^2)$,
i.e.\ $\lambda\ge 4BC\{\sqrt{\log d/n}+\delta\}$ (after $(a+b)^2\ge a^2+b^2$ and $2d\,d^{-2}\le 2$). The
factor $2BC$ that suffices for the deterministic rate is \emph{not} enough to place the realized
gradient inside $\mathfrak{G}(\lambda)$ with the claimed probability. We therefore state
Corollary~9.27 with the \emph{same} $\lambda=4BC\{\cdots\}$ as Corollary~9.26; the deterministic
constant $3\lambda/\kappa$ in the conclusion is unaffected.

\section{Discussion}\label{sec:disc}

\paragraph{The rate is untouched.} Every discrepancy above is a change of \emph{constant}. The
$\ell_2$, $\ell_1$, and $\ell_\infty$ rates remain
$\norm{\widehat\Delta}^2 = O\!\big(s\lambda^2/\kappa^2\big)$,
$\norm{\widehat\Delta}_1 = O\!\big(s\lambda/\kappa\big)$,
$\norm{\widehat\Delta}_\infty = O\!\big(\lambda/\kappa\big)$ with
$\lambda\asymp\sqrt{\log d/n}$, exactly as in the book. No exponent, no dependence on $(n,d,s,\kappa)$
changes.

\paragraph{The root cause is structural, not a book error.} Wainwright's $9$ is the correct
\emph{leading-order} constant in the regime $\tau_n^2\PsM^2/\kappa\to0$, which is how the theory is used
asymptotically (the slack $\tau_n^2$ shrinks like $\log d/n$). A formal proof, however, asserts a bound
that holds at \emph{every} finite $n$ satisfying the stated hypotheses, including those where
$\tau_n^2\PsM^2$ is a genuine (if small) fraction of $\kappa$. There the quadratic defect~\eqref{eq:defectquad}
cannot be discarded, and the honest, explicit constant is larger. This is an instance of the general
caveat recorded in the project charter: \emph{textbook constants are stated up to small factors; the
formalization states what is actually provable and documents the deviation in the declaration's
docstring.}

\paragraph{Could the constant be improved?} Yes, at a cost in generality. A bespoke $\thetast\in M$
argument that never introduces the approximation cross-term keeps $c_{\mathrm{def}}=16$ and the budget
$\kappa/4$ undivided, yielding $\norm{\widehat\Delta}\le 6\lambda\PsM/\kappa$, i.e.\ the constant $36$
with slack $\kappa/64$. We deliberately formalized the single \emph{general} Theorem~9.19 for
$M\subseteq\Mbar$ (the book-faithful scope) and obtained Corollary~9.20 by specialization, accepting the
inherited $144$ rather than maintaining a second, separately tuned proof of the special case. Sharper
constants would require either (i) a dedicated $\thetast\in M$ theorem, or (ii) a localized argument that
re-introduces the defect at the critical radius with a self-bounding step; neither changes the rate, so
we did not pursue them.

\paragraph{Bookkeeping.} All deviations are tagged in the Lean sources:
\texttt{Bound.lean} (the $144$/$32$ constants and the $\kappa/128$ slack, with the
$\kappa/4=\kappa/8+\kappa/16+\kappa/16$ split), \texttt{DualBound.lean} (the $\kappa/32$ slack of
Theorem~9.24, which \emph{does} match the book), and \texttt{GLMCorollaries.lean} (the
$\lambda=4BC\{\cdots\}$ reuse in Corollary~9.27). Each main theorem is \texttt{\#print axioms}--clean.

\begin{thebibliography}{9}
\bibitem{wainwright} M.~J.\ Wainwright, \emph{High-Dimensional Statistics: A Non-Asymptotic Viewpoint},
Cambridge University Press, 2019. Chapter~9 (Theorems 9.19, 9.24; Corollaries 9.20, 9.26, 9.27).
\bibitem{statlean} StatLean formalization, \texttt{StatLean/HighDimensionalStatistics/MEstimator/}
(Batch~6); milestone notes \texttt{notes/high\_dimensional\_statistics/m\_estimator/}.
\end{thebibliography}

\end{document}
